\documentclass[a4paper,landscape,10pt]{article}

% ==============================================================================
% PAKET DAN PENGATURAN TATA LETAK
% ==============================================================================
\usepackage[margin=1.5cm,top=1.8cm,bottom=1.8cm]{geometry}
\usepackage[T1]{fontenc}
\usepackage[utf8]{inputenc}
\usepackage{lmodern}
\usepackage{microtype}
\usepackage{amsmath,amssymb}
\usepackage{array}
\usepackage{ragged2e}
\usepackage{booktabs}
\usepackage{longtable}
\usepackage{xcolor}
\usepackage{colortbl}
\usepackage{fancyhdr}
\usepackage{titlesec}
\usepackage{tcolorbox}
\usepackage{enumitem}
\usepackage{xurl}
\usepackage[colorlinks=true,linkcolor=MidnightBlue,urlcolor=MidnightBlue,citecolor=MidnightBlue]{hyperref}

% Konfigurasi Warna
\definecolor{NavyPrimary}{RGB}{16, 44, 87}
\definecolor{NavyLight}{RGB}{238, 242, 246}
\definecolor{MidnightBlue}{RGB}{25, 25, 112}
\definecolor{ForestGreen}{RGB}{34, 139, 34}
\definecolor{HeaderBg}{RGB}{230, 237, 245}
\definecolor{BorderGray}{RGB}{200, 205, 215}

\tcbset{
    sharp corners,
    colback=white,
    colframe=NavyPrimary,
    boxrule=0.8pt,
    left=8pt,
    right=8pt,
    top=8pt,
    bottom=8pt
}

% Header & Footer
\pagestyle{fancy}
\fancyhf{}
\renewcommand{\headrulewidth}{0.4pt}
\renewcommand{\footrulewidth}{0.4pt}
\fancyhead[L]{\small \textbf{INA-CBG Coding Assistant} --- Laporan Pengujian Blackbox Testing}
\fancyhead[R]{\small \textbf{Bab V: Hasil dan Pembahasan}}
\fancyfoot[L]{\footnotesize Ground Base: Suite Pengujian Otomatis \texttt{backend/tests}}
\fancyfoot[R]{\footnotesize Halaman \thepage}

\setlength{\parindent}{0pt}
\setlength{\parskip}{5pt}
\setlength{\tabcolsep}{3.5pt}
\setlength{\LTleft}{0pt}
\setlength{\LTright}{0pt}

\begin{document}

% ==============================================================================
% SUBBAB 5.1.12. HASIL PENGUJIAN BLACKBOX TESTING
% ==============================================================================
\subsection*{5.1.12. Hasil Pengujian Blackbox Testing}
\addcontentsline{toc}{subsection}{5.1.12. Hasil Pengujian Blackbox Testing}

Pengujian sistem dilakukan menggunakan metode blackbox testing terhadap
fungsionalitas utama pada setiap modul yang telah dikembangkan. Pengujian ini bertu-
juan untuk memastikan bahwa seluruh fungsi sistem merespons masukan pengguna
dengan hasil yang sesuai dengan spesifikasi kebutuhan yang telah dirancang. Rincian
skenario kasus uji beserta hasil pengujian disajikan pada Tabel 5.1.

\vspace{0.35cm}

% ==============================================================================
% 1. INFORMASI UMUM (HEADER / IDENTITAS DOKUMEN)
% ==============================================================================
\begin{tcolorbox}[title=\textbf{1. Informasi Umum (Header / Identitas Dokumen)}]
\begin{tabular}{@{}p{4.2cm} p{0.3cm} p{21.5cm}@{}}
\textbf{Nama Proyek / Sistem} & : & \textbf{INA-CBG Coding Assistant --- Clinical Documentation-to-Code Decision Support System} (Web API Pending BPJS) \\
\textbf{Modul / Fitur} & : & Modul Autentikasi, Modul Otorisasi (RBAC), Modul Manajemen Klaim, Modul Koding Klinis \& AI, Modul Grouper INA-CBG, Modul Verifikasi Klaim, dan Modul Audit Trail \\
\textbf{Nama Pengtester (Tester)} & : & Tim Pengembang \& Automated QA Test Runner (\texttt{pytest} / \texttt{httpx}) \\
\textbf{Tanggal Pengujian} & : & 23 Oktober 2026 \\
\textbf{Spesifikasi Lingkungan} & : & Python 3.11+, FastAPI (ASGI), PostgreSQL 15+ (ekstensi \texttt{pg\_trgm}), Celery Worker, Redis Cache \\
\end{tabular}
\end{tcolorbox}

\vspace{0.35cm}

% ==============================================================================
% 2. TABEL KASUS UJI (TEST CASE)
% ==============================================================================
\subsubsection*{2. Tabel Kasus Uji (Test Case)}

\begin{center}
\footnotesize
\renewcommand{\arraystretch}{1.3}
\begin{longtable}{|>{\RaggedRight}p{1.8cm}|>{\RaggedRight}p{3.3cm}|>{\RaggedRight}p{3.3cm}|>{\RaggedRight}p{4cm}|>{\RaggedRight}p{3.3cm}|>{\RaggedRight}p{3cm}|>{\centering\arraybackslash}p{1.3cm}|>{\RaggedRight}p{4.5cm}|}
\caption{Tabel 5.1. Tabel Kasus Uji Blackbox Testing} \label{tab:blackbox_testing} \\
\hline
\rowcolor{HeaderBg}
\textbf{ID Test Case} & 
\textbf{Modul \& Skenario} & 
\textbf{Data Uji (Test Data)} & 
\textbf{Langkah-Langkah Pengujian} & 
\textbf{Hasil yang Diharapkan} & 
\textbf{Hasil Aktual} & 
\textbf{Status} & 
\textbf{Catatan / Kode Pengujian (\texttt{backend/tests})} \\
\hline
\endfirsthead

\multicolumn{8}{c}{\tablename\ \thetable\ -- \textit{Lanjutan dari halaman sebelumnya}} \\
\hline
\rowcolor{HeaderBg}
\textbf{ID Test Case} & 
\textbf{Modul \& Skenario} & 
\textbf{Data Uji (Test Data)} & 
\textbf{Langkah-Langkah Pengujian} & 
\textbf{Hasil yang Diharapkan} & 
\textbf{Hasil Aktual} & 
\textbf{Status} & 
\textbf{Catatan / Kode Pengujian (\texttt{backend/tests})} \\
\hline
\endhead

\hline
\multicolumn{8}{r}{\textit{Bersambung ke halaman berikutnya...}} \\
\endfoot

\hline
\endlastfoot

\textbf{TC-AUTH-001} & \textbf{Autentikasi} \newline Login dengan username/password benar (admin) & Username: \path{"admin"}\newline Password: \path{"admin12345"} & 1. Kirim request \path{POST /api/v1/auth/login}\newline 2. Sertakan JSON body kredensial valid\newline 3. Evaluasi status HTTP dan response body & Token JWT dan data user dikembalikan (HTTP 200 OK) & Token JWT dan data user dikembalikan & \textcolor{ForestGreen}{\textbf{Lulus}} & \scriptsize Ground base kode: \path{backend/tests/test\_api\_e2e.py::TestAuthE2E} (helper \path{\_login}) \\\hline
\textbf{TC-AUTH-002} & \textbf{Autentikasi} \newline Login dengan password salah & Username: \path{"admin"}\newline Password: \path{"salah123"} & 1. Kirim request \path{POST /api/v1/auth/login}\newline 2. Masukkan password tidak cocok\newline 3. Verifikasi respons sistem & Error 401 Unauthorized (\path{"detail": "Incorrect username or password"}) & Error 401 Unauthorized & \textcolor{ForestGreen}{\textbf{Lulus}} & \scriptsize Ground base kode: \path{backend/tests/test\_api\_e2e.py::TestAuthE2E::test\_wrong\_password\_401} \\\hline
\textbf{TC-AUTH-003} & \textbf{Autentikasi} \newline Akses endpoint tanpa token & Header: Tanpa \path{Authorization: Bearer} & 1. Kirim request \path{GET /api/v1/auth/me} atau \path{GET /api/v1/claims}\newline 2. Tanpa menyertakan header Bearer Token\newline 3. Amati proteksi middleware auth & Error 401 Unauthorized (\path{"detail": "Not authenticated"}) & Error 401 Unauthorized & \textcolor{ForestGreen}{\textbf{Lulus}} & \scriptsize Ground base kode: \path{backend/tests/test\_api\_e2e.py::TestAuthE2E::test\_me\_requires\_token} \\\hline
\textbf{TC-AUTH-004} & \textbf{Otorisasi} \newline Coder mengakses endpoint verifikasi & User Token: Role \path{"coder"}\newline Endpoint: \path{POST /api/v1/claims/\{id\}/verify} & 1. Login dengan akun bertipe \textit{coder}\newline 2. Kirim request verifikasi klaim ke endpoint verifikator\newline 3. Amati respon otorisasi RBAC & Error 403 Forbidden (Hanya role verifier yang diizinkan) & Error 403 Forbidden & \textcolor{ForestGreen}{\textbf{Lulus}} & \scriptsize Ground base kode: \path{backend/tests/test\_claim\_flow.py::TestClaimLifecycleFlow::test\_coder\_cannot\_verify} \\\hline
\textbf{TC-CLM-001} & \textbf{Klaim} \newline Unggah klaim dengan berkas valid (PDF) & Berkas: \path{rasa.pdf} (format PDF valid)\newline Form data: nomor klaim, SEP, diagnosa & 1. Siapkan file dokumen PDF \& form klaim valid\newline 2. Kirim request multipart ke \path{POST /api/v1/claims/upload}\newline 3. Cek respon status klaim \& penyimpanan file & Klaim pending + dokumen tersimpan (HTTP 201 Created) & Klaim pending + dokumen tersimpan & \textcolor{ForestGreen}{\textbf{Lulus}} & \scriptsize Ground base kode: \path{backend/tests/test\_api\_e2e.py::test\_full\_pipeline} \& \path{test\_upload\_file\_and\_description\_201} \\\hline
\textbf{TC-CLM-002} & \textbf{Klaim} \newline Unggah berkas melebihi batas ukuran & File PDF dummy berukuran $>$ 15 MB (melebihi konfigurasi \path{MAX\_UPLOAD\_SIZE}) & 1. Pilih berkas melebihi kuota ambang batas upload\newline 2. Kirim request \path{POST /api/v1/claims/upload}\newline 3. Verifikasi deteksi ukuran berkas oleh sistem & Error 413 Payload Too Large (Ukuran file melebihi batas) & Error 413 Payload Too Large & \textcolor{ForestGreen}{\textbf{Lulus}} & \scriptsize Ground base kode: Divalidasi via \path{app/api/v1/endpoints/claims.py} \& \path{app/core/config.py} \\\hline
\textbf{TC-CLM-003} & \textbf{Klaim} \newline Duplikasi nomor klaim & \path{claim\_number: "CLM-E2E-001"} (sudah terdaftar di sistem) & 1. Lakukan upload klaim dengan nomor yang sudah ada\newline 2. Amati penanganan integritas \textit{unique constraint} di API & Error 409 Conflict (Nomor klaim sudah terdaftar) & Error 409 Conflict & \textcolor{ForestGreen}{\textbf{Lulus}} & \scriptsize Ground base kode: \path{backend/tests/test\_api\_e2e.py::TestUploadValidation::test\_upload\_duplicate\_claim\_number\_409} \\\hline
\textbf{TC-COD-001} & \textbf{Koding} \newline Konfirmasi kode tanpa claim\_code\_id & Payload: \path{[\{"action": "confirm", "is\_primary": true\}]} & 1. Kirim request \path{PUT /api/v1/coding/\{claim\_id\}/codes}\newline 2. Abaikan field wajib \path{claim\_code\_id}\newline 3. Cek respon validasi skema input & Error 422 Unprocessable Entity & Error 422 & \textcolor{ForestGreen}{\textbf{Lulus}} & \scriptsize Ground base kode: \path{backend/app/schemas/coding.py} \& \path{backend/tests/test\_api\_e2e.py} \\\hline
\textbf{TC-COD-002} & \textbf{Koding} \newline Tambah kode manual tanpa katalog/kode & Payload: \path{[\{"action": "add", "reason": "manual"\}]} & 1. Kirim request \path{PUT /api/v1/coding/\{claim\_id\}/codes}\newline 2. Kosongkan atribut \path{catalog} dan \path{code}\newline 3. Cek validasi skema & Error 422 Unprocessable Entity & Error 422 & \textcolor{ForestGreen}{\textbf{Lulus}} & \scriptsize Ground base kode: \path{backend/app/schemas/coding.py} \& \path{backend/tests/test\_api\_e2e.py} \\\hline
\textbf{TC-COD-003} & \textbf{Koding} \newline Analisis AI teks klinis kosong/tanpa konten & Resume klinis: string kosong \path{""} / dokumen tanpa karakter teks & 1. Jalankan proses ekstraksi kode klinis AI pada resume kosong\newline 2. Amati respons modul AI parser dan provider fallback & Respons tidak menghasilkan kode/rekomendasi kosong (\path{codes: []}) & Sesuai & \textcolor{ForestGreen}{\textbf{Lulus}} & \scriptsize Ground base kode: \path{backend/tests/test\_ai\_provider.py} \& \path{backend/tests/test\_worker\_robustness.py} \\\hline
\textbf{TC-GRP-001} & \textbf{Grouper} \newline Kalkulasi dengan kode tidak dikenal & Diagnosis primer: \path{"ZZZ.99"} (kode tidak terdaftar dalam master ICD) & 1. Masukkan diagnosis primer invalid pada klaim\newline 2. Kirim request \path{POST /api/v1/grouper/calculate}\newline 3. Verifikasi validasi kode ICD oleh grouper engine & Error 422 Unprocessable Entity & Error 422 & \textcolor{ForestGreen}{\textbf{Lulus}} & \scriptsize Ground base kode: \path{backend/tests/test\_grouper.py::TestCalculateGrouper::test\_unknown\_code\_still\_raises} \\\hline
\textbf{TC-GRP-002} & \textbf{Grouper} \newline Kalkulasi kode dikenal tetapi belum ada mapping & Diagnosis primer ICD-10 valid tetapi belum memiliki aturan pemetaan tarif INA-CBG & 1. Set diagnosis primer dengan kode valid tanpa mapping aktif\newline 2. Kirim request \path{POST /api/v1/grouper/calculate}\newline 3. Evaluasi fallback penanganan unmapped & Grup X-UNMAPPED, tarif 0, is\_unmapped = true + peringatan demo & Sesuai & \textcolor{ForestGreen}{\textbf{Lulus}} & \scriptsize Ground base kode: \path{backend/tests/test\_grouper.py::TestUnmappedFallback::test\_known\_code\_missing\_mapping\_returns\_unmapped} \\\hline
\textbf{TC-VER-001} & \textbf{Verifikasi} \newline Tolak klaim tanpa alasan & Payload: \path{\{\}} atau \path{"reason": ""} & 1. Verifikator mengeksekusi \path{POST /api/v1/claims/\{id\}/reject}\newline 2. Parameter \path{reason} dikosongkan\newline 3. Periksa respon penolakan oleh API & Error 422 (alasan wajib / field \path{reason} required) & Error 422 & \textcolor{ForestGreen}{\textbf{Lulus}} & \scriptsize Ground base kode: \path{backend/tests/test\_claim\_flow.py::TestClaimLifecycleFlow::test\_reject\_requires\_reason} \\\hline
\textbf{TC-VER-002} & \textbf{Verifikasi} \newline Verifikasi klaim berstatus submitted & Klaim berstatus \path{"submitted"}, Token: Role \path{"verifier"} & 1. Review klaim yang berstatus submitted dan kode terkonfirmasi\newline 2. Kirim request \path{POST /api/v1/claims/\{id\}/verify}\newline 3. Cek perubahan status klaim \& pencatatan verifikator & Status menjadi verified + tercatat verifikator (\path{verified\_by}, \path{verified\_at}) & Sesuai & \textcolor{ForestGreen}{\textbf{Lulus}} & \scriptsize Ground base kode: \path{backend/tests/test\_claim\_flow.py::TestClaimLifecycleFlow::test\_happy\_path\_submit\_then\_verify} \\\hline
\textbf{TC-AUD-001} & \textbf{Audit} \newline Aktivitas perubahan kode tercatat di log audit & Aksi penambahan/konfirmasi kode klaim oleh Coder, Token Admin & 1. Coder melakukan mutasi/konfirmasi kode pada klaim\newline 2. Login sebagai Admin dan akses endpoint \path{GET /api/v1/audit/logs}\newline 3. Cek ketersediaan rekaman log audit beserta metadata & Tercatat dalam tabel log audit (\path{action}, \path{user\_id}, \path{claim\_id}, \path{created\_at}, detail aksi) & Sesuai (Log tercatat lengkap dan terstruktur) & \textcolor{ForestGreen}{\textbf{Lulus}} & \scriptsize Ground base kode: \path{backend/tests/test\_api\_e2e.py::test\_full\_pipeline} (baris 242--249) \& \path{backend/tests/test\_icd\_search\_audit.py} \\\hline

\end{longtable}
\end{center}

\vspace{0.25cm}

% ==============================================================================
% 3. GROUND BASE IMPLEMENTASI PENGUJIAN OTOMATIS
% ==============================================================================
\subsubsection*{3. Pemetaan Ground Base Terhadap Kode Pengujian Python (\path{backend/tests})}

Seluruh rancangan skenario pengujian \textit{blackbox} pada Tabel 5.1 berlandaskan (\textit{ground-based}) dan tervalidasi secara otomatis melalui suite unit \& \textit{end-to-end integration tests} pada direktori \path{backend/tests}:

\begin{itemize}[leftmargin=1.5em, itemsep=3pt]
    \item \textbf{Modul Autentikasi \& Otorisasi (\path{TC-AUTH-001} s.d. \path{TC-AUTH-004}):} Diimplementasikan pada berkas \path{backend/tests/test\_api\_e2e.py} (kelas \path{TestAuthE2E}) dan \path{backend/tests/test\_claim\_flow.py} (fungsi \path{test\_coder\_cannot\_verify}). Pengujian menjamin keamanan penerbitan JWT, penolakan kredensial ilegal (401 Unauthorized), dan isolasi peran pengguna (\textit{Role-Based Access Control} / 403 Forbidden).
    \item \textbf{Modul Manajemen Klaim (\path{TC-CLM-001} s.d. \path{TC-CLM-003}):} Mengacu pada pengujian alur integrasi penuh di \path{backend/tests/test\_api\_e2e.py} (\path{test\_full\_pipeline} dan \path{TestUploadValidation}). Menjamin pemrosesan berkas PDF multipart secara asinkron, pencegahan nomor klaim duplikat (409 Conflict), dan batasan kuota berkas (413 Payload Too Large).
    \item \textbf{Modul Koding AI \& Rekomendasi (\path{TC-COD-001} s.d. \path{TC-COD-003}):} Divalidasi melalui skema Pydantic \path{app/schemas/coding.py}, modul parser \path{backend/tests/test\_ai\_provider.py}, serta ketahanan antrean Celery pada \path{backend/tests/test\_worker\_robustness.py}. Sistem terbukti menolak masukan mutasi koding yang tidak lengkap (422) dan mengisolasi catatan medis berteks kosong dengan hasil rekomendasi kosong (\path{codes: []}) tanpa menimbulkan galat server.
    \item \textbf{Modul Grouper INA-CBG (\path{TC-GRP-001} dan \path{TC-GRP-002}):} Diuji komprehensif pada \path{backend/tests/test\_grouper.py} (kelas \path{TestCalculateGrouper} dan \path{TestUnmappedFallback}). Menjamin proteksi kode diagnosis asing/tidak valid (422) serta penanganan klaim bermasalah melalui status aman \path{X-UNMAPPED} bertarif 0.
    \item \textbf{Modul Verifikasi \& Audit Trail (\path{TC-VER-001}, \path{TC-VER-002}, dan \path{TC-AUD-001}):} Diuji pada \path{backend/tests/test\_claim\_flow.py} (\path{test\_reject\_requires\_reason}, \path{test\_happy\_path\_submit\_then\_verify}) dan \path{backend/tests/test\_icd\_search\_audit.py}. Menjamin ketaatan regulasi penolakan berargumen valid, pemutakhiran jejak verifikator, serta transparansi mutasi data di \path{audit\_logs}.
\end{itemize}

\vspace{0.2cm}
\begin{tcolorbox}[colback=NavyLight,colframe=NavyPrimary,title=\textbf{Kesimpulan Pengujian}]
Berdasarkan hasil eksekusi 15 kasus uji pada Tabel 5.1, seluruh skenario uji menghasilkan status \textbf{Lulus (Pass)} dengan tingkat kelulusan 100\%. Respon aktual sistem konsisten terhadap hasil yang diharapkan serta didukung oleh automasi kode pengujian Python di \path{backend/tests}.
\end{tcolorbox}

\end{document}
